\documentclass[tikz,border=10pt]{standalone}
\input{../../../docs/tikz/dag_preamble.tex}

\begin{document}
\begin{tikzpicture}[x=1cm,y=1cm,every node/.append style={outer sep=4pt}]
\dagPaperMetadata
  {Capacity Constraints Make Admissions Processes Less Predictable}
  {Evan Dong, Nikhil Garg, and Sarah Dean}
  {AAAI-26}
  {2026}
  {https://ojs.aaai.org/index.php/AAAI/article/view/41179}

\dagPaperLegendRightOfMetadata{
\node[dag_result,dag_template_legend] (legExact) at (0,0) {source-shaped\\formalization};
\node[dag_template_note,dag_template_legend] (legNote) at (4.6,0) {source-model\\clarification};
}
\daglegend{(legExact)(legNote)}{Legend}

\begin{dagPaperBody}
\node[dag_result] (CoreModel) at (0,0) {\textbf{Finite Choice Model}\\feasibility, capacity, distance,\\instability and variability};
\node[dag_result] (TheoremOne) at (6.2,0) {\textbf{Theorem 1}\\no-zero nondegeneracy,\\substitutability and tight families};
\node[dag_result] (MLFormulas) at (12.4,0) {\textbf{ML Representation}\\choice labels, fixed threshold,\\rank threshold and score embedding};

\node[dag_result] (Sequential) at (0,-4.1) {\textbf{Sequential Composition}\\feasibility, substitutability,\\one-instability and additive bounds};
\node[dag_result] (TheoremTwo) at (6.2,-4.1) {\textbf{Theorem 2}\\realized range $1\leq m\leq n$ and\\single-order characterization};
\node[dag_result] (Programs) at (12.4,-4.1) {\textbf{Proposition 2 Programs}\\abstract queue procedures; exact\\variability $1$, $2$, $3$, and $6$};

\node[dag_result] (ChoiceProps) at (0,-8.2) {\textbf{Choice-Property Calculus}\\monotonicity, consistency,\\independence and equivalences};
\node[dag_result] (Parity) at (6.2,-8.2) {\textbf{Parity Characterization}\\positive even distance iff inconsistency;\\no consistent positive even tightness};
\node[dag_result] (AppendRemove) at (12.4,-8.2) {\textbf{Append/Remove Variability}\\borderline--waitlisted equality and\\cross-pool removable-set equality};

\node[dag_result] (LAPModel) at (0,-12.3) {\textbf{Source LAP Model}\\real weights, no ties, feasibility,\\capacity and objective optimality};
\node[dag_result] (LAPOrder) at (6.2,-12.3) {\textbf{LAP Ordering}\\assigned applicants outrank rejected;\\higher applicants are selected};
\node[dag_result] (LAPVariation) at (12.4,-12.3) {\textbf{LAP Predictability}\\one-instability and variability bounded\\by distinct induced slot orders};

\node[dag_template_note,text width=18.2cm] (StatusNote) at (6.2,-16.4) {\textbf{Status: formalized.}  All 64 selected source-to-Spec claims are covered.  Proposition 2 is formalized as the paper's abstract one-, two-, three-, and six-queue procedures, not as an empirical roster or simulator.  Fixed generic tie-breaking is a source-model clarification that makes the operational selector single-valued; it is not an economic restriction or a caveat.};

\draw[dag_arrow] (CoreModel) -- (TheoremOne);
\draw[dag_arrow] (TheoremOne) -- (MLFormulas);
\draw[dag_arrow] (CoreModel) -- (Sequential);
\draw[dag_arrow] (Sequential) -- (TheoremTwo);
\draw[dag_arrow] (TheoremTwo) -- (Programs);
\draw[dag_arrow] (Sequential) -- (ChoiceProps);
\draw[dag_arrow] (ChoiceProps) -- (Parity);
\draw[dag_arrow] (Parity) -- (AppendRemove);
\draw[dag_arrow] (ChoiceProps) -- (LAPModel);
\draw[dag_arrow] (LAPModel) -- (LAPOrder);
\draw[dag_arrow] (LAPOrder) -- (LAPVariation);
\end{dagPaperBody}
\end{tikzpicture}
\end{document}
